\section{Simulation construction and complete results}\label{supp:simulation}
\newcommand{\JointActiveCoverage}{0.964}
\newcommand{\JointActiveMCSE}{0.004}
\newcommand{\JointActiveAblationCoverage}{0.935}
\newcommand{\JointActiveCoverageDifference}{0.028}
\newcommand{\JointActiveDifferenceMCSE}{0.004}
\newcommand{\JointActiveLength}{0.275}
\newcommand{\JointActiveAblationLength}{0.259}
\newcommand{\JointNullCoverage}{0.968}
\newcommand{\JointNullMCSE}{0.003}
\newcommand{\JointNullAblationCoverage}{0.956}
\newcommand{\JointNullCoverageDifference}{0.012}
\newcommand{\JointNullDifferenceMCSE}{0.004}
\newcommand{\JointNullLength}{0.315}
\newcommand{\JointNullAblationLength}{0.286}
\newcommand{\JointDatasets}{1440}
\newcommand{\JointIntervals}{5760}
\newcommand{\JointAblationShiftCount}{306}
\newcommand{\JointAblationMaxShift}{0.350}
\newcommand{\JointWeakCoverage}{0.958}
\newcommand{\JointWeakAblationCoverage}{0.803}
\newcommand{\JointWeakDifference}{0.156}
\newcommand{\JointWeakDifferenceMCSE}{0.022}
\newcommand{\JointRuntimeFiveHundred}{0.966}
\newcommand{\JointRuntimeTenThousand}{9.976}
\newcommand{\JointSingleMinimumOrder}{1.003}
\newcommand{\JointSingleMaximumOrder}{1.014}
\newcommand{\JointPairMinimumOrder}{1.998}
\newcommand{\JointPairMaximumOrder}{2.000}
\newcommand{\JointModuleWorkflows}{52}
\newcommand{\JointModulePairedFits}{6760}
\newcommand{\JointPlasmaALowerEstimate}{-0.133}
\newcommand{\JointPlasmaALowerScoreInterval}{[-0.213, -0.054]}
\newcommand{\JointPlasmaALowerWorkflowInterval}{[-0.222, -0.045]}
\newcommand{\JointPlasmaAMedianEstimate}{-0.091}
\newcommand{\JointPlasmaAMedianScoreInterval}{[-0.163, -0.019]}
\newcommand{\JointPlasmaAMedianWorkflowInterval}{[-0.168, -0.013]}
\newcommand{\JointPlasmaAUpperEstimate}{-0.057}
\newcommand{\JointPlasmaAUpperScoreInterval}{[-0.123, 0.009]}
\newcommand{\JointPlasmaAUpperWorkflowInterval}{[-0.125, 0.011]}
\newcommand{\JointPlasmaBLowerEstimate}{-0.068}
\newcommand{\JointPlasmaBLowerScoreInterval}{[-0.144, 0.009]}
\newcommand{\JointPlasmaBLowerWorkflowInterval}{[-0.151, 0.015]}
\newcommand{\JointPlasmaBMedianEstimate}{-0.065}
\newcommand{\JointPlasmaBMedianScoreInterval}{[-0.122, -0.008]}
\newcommand{\JointPlasmaBMedianWorkflowInterval}{[-0.126, -0.003]}
\newcommand{\JointPlasmaBUpperEstimate}{-0.036}
\newcommand{\JointPlasmaBUpperScoreInterval}{[-0.089, 0.018]}
\newcommand{\JointPlasmaBUpperWorkflowInterval}{[-0.092, 0.021]}
\newcommand{\JointGeneRoles}{360}
\newcommand{\JointGeneIntervals}{1080}
\newcommand{\JointLowerGeneA}{PTTG1}
\newcommand{\JointLowerGeneCountA}{11}
\newcommand{\JointLowerGeneB}{PTTG1}
\newcommand{\JointLowerGeneCountB}{6}

\newcommand{\RegisteredJointSettings}{71}
\newcommand{\RegisteredJointDatasets}{45500}
\newcommand{\RegisteredJointFailures}{6}
\newcommand{\RegisteredJointAvailable}{45494}
\newcommand{\RegisteredFullPredictorSettings}{66}
\newcommand{\RegisteredFullActiveCoverage}{0.971}
\newcommand{\RegisteredFullActiveMCSE}{0.001}
\newcommand{\RegisteredFullNullCoverage}{0.977}
\newcommand{\RegisteredFullNullMCSE}{0.001}
\newcommand{\VarianceFilterActiveCoverage}{0.661}
\newcommand{\VarianceFilterNullCoverage}{0.694}
\newcommand{\FullDesignActiveCoverage}{0.977}
\newcommand{\FullDesignNullCoverage}{0.989}
\newcommand{\VarianceFilterSecondShifted}{463}
\newcommand{\VarianceFilterEleventhShifted}{472}


\subsection{Registered statistical configurations}\label{supp:registered-study}

The registered study comprises \RegisteredJointSettings{} settings
and \RegisteredJointDatasets{} independently generated datasets.
Table~\ref{tab:registered-study-design} gives their sample sizes,
dimensions and replication counts. Each dataset requests coordinates
$1$, $2$, $s+1$ and $s+2$, where $s$ is the number of nonzero
generating coefficients. The first two coordinates assess active
effects and the latter two assess null effects. The target is the
complete unsmoothed profile parameter for the full generating design.

The Gaussian scaling settings A01 through A12 vary the cluster count
at each dimension and sparsity level. The error settings B01 through
B19 use standard normal, $t_3/\sqrt{3}$, unit-variance Laplace and
$(\chi_3^2-3)/\sqrt{6}$ innovations, together with a
$0.9N(0,1)+0.1N(0,25)$ mixture and a unit-variance asymmetric Laplace
distribution. Each innovation is centred at its specified quantile.
Predictor-dependent scale multiplies the centred innovation by
$\exp\{0.35\operatorname{clip}(X_1,-3,3)\}$. The nonmedian
scale settings B16 and B19 use the independently calculated profile
targets described below.

Cluster heterogeneity settings C01 through C16 change the random
intercept variance and include correlated intercepts and time slopes.
They also examine non-Gaussian cluster effects and departures from
the fitted nuisance model. Gaussian-copula innovations have
within-cluster AR(1) correlation $0.4$ and preserve their declared
marginal quantiles. The non-informative imbalance setting samples
$m\in\{2,\ldots,12\}$ with probability proportional to
$0.35(1-0.35)^{m-2}$. Profile target calculations apply to the omitted
slope, random-scale and nonlinear nuisance settings, and to the
settings that link covariates or cluster size to cluster effects.

The design settings D01 through D06 vary correlation and the
population inverse-information structure. The working-penalty
settings retain fixed positive values of $\Lambda$, and the bandwidth
sensitivities use multiples $0.75$ and $1.25$ of $n^{-3/10}$.
E07 and E08 multiply the joint learning budget by $0.5$ and $2$.
The signal settings use coefficient magnitudes $0.40$ and $1.10$.
The complete generating parameters are retained in the accompanying
configuration and source tables.

In F01 through F03, independent screening uses half of the subjects
and passes its predictor pool to the estimation sample. F04 uses
same-sample CA-IQR-SIS. F05 uses outcome-free variance filtering,
while F06 passes the full predictor set to coordinate learning.
Each screening rule adds the requested coordinates when needed.
The coverage calculations retain the full-design targets across
these input rules.

\begingroup\small

\begin{longtable}[t]{lrrrrrrl}
\caption{\label{tab:registered-study-design}Registered statistical settings and dataset counts}\\
\toprule
Setting & $n$ & $p$ & $s$ & $\tau$ & $q$ & Replications & Input\\
\midrule
\endfirsthead
\caption[]{Registered statistical settings and dataset counts \textit{(continued)}}\\
\toprule
Setting & $n$ & $p$ & $s$ & $\tau$ & $q$ & Replications & Input\\
\midrule
\endhead

\endfoot
\bottomrule
\endlastfoot
A01 & 100 & 500 & 5 & 0.50 & 1 & 1000 & Full predictor set\\
A02 & 200 & 500 & 5 & 0.50 & 1 & 1000 & Full predictor set\\
A03 & 400 & 500 & 5 & 0.50 & 1 & 1000 & Full predictor set\\
A04 & 100 & 1000 & 5 & 0.50 & 1 & 1000 & Full predictor set\\
A05 & 200 & 1000 & 5 & 0.50 & 1 & 1000 & Full predictor set\\
A06 & 400 & 1000 & 5 & 0.50 & 1 & 1000 & Full predictor set\\
A07 & 100 & 2000 & 5 & 0.50 & 1 & 500 & Full predictor set\\
A08 & 200 & 2000 & 5 & 0.50 & 1 & 500 & Full predictor set\\
A09 & 400 & 2000 & 5 & 0.50 & 1 & 500 & Full predictor set\\
A10 & 100 & 1000 & 10 & 0.50 & 1 & 1000 & Full predictor set\\
A11 & 200 & 1000 & 10 & 0.50 & 1 & 1000 & Full predictor set\\
A12 & 400 & 1000 & 10 & 0.50 & 1 & 1000 & Full predictor set\\
B01 & 200 & 1000 & 8 & 0.25 & 1 & 1000 & Full predictor set\\
B02 & 200 & 1000 & 8 & 0.50 & 1 & 1000 & Full predictor set\\
B03 & 200 & 1000 & 8 & 0.75 & 1 & 1000 & Full predictor set\\
B04 & 200 & 1000 & 8 & 0.25 & 1 & 1000 & Full predictor set\\
B05 & 200 & 1000 & 8 & 0.50 & 1 & 1000 & Full predictor set\\
B06 & 200 & 1000 & 8 & 0.75 & 1 & 1000 & Full predictor set\\
B07 & 200 & 1000 & 8 & 0.25 & 1 & 500 & Full predictor set\\
B08 & 200 & 1000 & 8 & 0.75 & 1 & 500 & Full predictor set\\
B09 & 200 & 1000 & 8 & 0.25 & 1 & 500 & Full predictor set\\
B10 & 200 & 1000 & 8 & 0.75 & 1 & 500 & Full predictor set\\
B11 & 200 & 1000 & 8 & 0.25 & 1 & 1000 & Full predictor set\\
B12 & 200 & 1000 & 8 & 0.50 & 1 & 1000 & Full predictor set\\
B13 & 200 & 1000 & 8 & 0.75 & 1 & 1000 & Full predictor set\\
B14 & 200 & 1000 & 8 & 0.25 & 1 & 500 & Full predictor set\\
B15 & 200 & 1000 & 8 & 0.75 & 1 & 500 & Full predictor set\\
B16 & 200 & 1000 & 8 & 0.25 & 1 & 500 & Full predictor set\\
B17 & 200 & 1000 & 8 & 0.50 & 1 & 500 & Full predictor set\\
B18 & 200 & 1000 & 8 & 0.50 & 1 & 500 & Full predictor set\\
B19 & 200 & 1000 & 8 & 0.75 & 1 & 500 & Full predictor set\\
C01 & 200 & 1000 & 8 & 0.50 & 1 & 500 & Full predictor set\\
C02 & 200 & 1000 & 8 & 0.50 & 1 & 1000 & Full predictor set\\
C03 & 200 & 1000 & 8 & 0.50 & 2 & 1000 & Full predictor set\\
C04 & 200 & 1000 & 8 & 0.50 & 2 & 500 & Full predictor set\\
C05 & 200 & 1000 & 8 & 0.50 & 2 & 500 & Full predictor set\\
C06 & 200 & 300 & 8 & 0.50 & 2 & 500 & Full predictor set\\
C07 & 200 & 300 & 8 & 0.25 & 1 & 500 & Full predictor set\\
C08 & 300 & 1000 & 8 & 0.50 & 1 & 500 & Full predictor set\\
C09 & 150 & 1000 & 8 & 0.50 & 2 & 500 & Full predictor set\\
C10 & 200 & 1000 & 8 & 0.50 & 1 & 500 & Full predictor set\\
C11 & 200 & 1000 & 8 & 0.50 & 1 & 500 & Full predictor set\\
C12 & 200 & 1000 & 8 & 0.25 & 1 & 500 & Full predictor set\\
C13 & 200 & 1000 & 8 & 0.50 & 1 & 500 & Full predictor set\\
C14 & 200 & 300 & 8 & 0.50 & 1 & 500 & Full predictor set\\
C15 & 200 & 300 & 8 & 0.50 & 1 & 500 & Full predictor set\\
C16 & 200 & 300 & 8 & 0.50 & 1 & 500 & Full predictor set\\
D01 & 200 & 1000 & 8 & 0.50 & 1 & 500 & Full predictor set\\
D02 & 200 & 1000 & 8 & 0.50 & 1 & 500 & Full predictor set\\
D03 & 200 & 1000 & 8 & 0.50 & 1 & 500 & Full predictor set\\
D04 & 200 & 1000 & 8 & 0.50 & 1 & 500 & Full predictor set\\
D05 & 200 & 1000 & 8 & 0.50 & 1 & 500 & Full predictor set\\
D06 & 200 & 1000 & 8 & 0.50 & 1 & 500 & Full predictor set\\
E01 & 200 & 1000 & 8 & 0.50 & 1 & 500 & Full predictor set\\
E02 & 200 & 1000 & 8 & 0.50 & 1 & 500 & Full predictor set\\
E03 & 200 & 1000 & 8 & 0.50 & 1 & 500 & Full predictor set\\
E04 & 200 & 1000 & 8 & 0.50 & 1 & 500 & Full predictor set\\
E05 & 200 & 1000 & 8 & 0.50 & 1 & 500 & Full predictor set\\
E06 & 200 & 1000 & 8 & 0.50 & 1 & 500 & Full predictor set\\
E07 & 200 & 1000 & 8 & 0.50 & 1 & 500 & Full predictor set\\
E08 & 200 & 1000 & 8 & 0.50 & 1 & 500 & Full predictor set\\
F01 & 200 & 5000 & 10 & 0.50 & 1 & 500 & Independent screening\\
F02 & 400 & 10000 & 10 & 0.50 & 1 & 500 & Independent screening\\
F03 & 300 & 5000 & 10 & 0.25 & 2 & 500 & Independent screening\\
F04 & 200 & 5000 & 10 & 0.50 & 1 & 500 & Same sample screening\\
F05 & 200 & 5000 & 10 & 0.50 & 1 & 500 & Variance filtering\\
F06 & 200 & 5000 & 10 & 0.50 & 1 & 500 & Full predictor set\\
H01 & 200 & 1000 & 8 & 0.50 & 1 & 500 & Full predictor set\\
H02 & 200 & 1000 & 8 & 0.50 & 1 & 500 & Full predictor set\\
H03 & 200 & 1000 & 8 & 0.25 & 1 & 500 & Full predictor set\\
H04 & 200 & 1000 & 8 & 0.25 & 1 & 500 & Full predictor set\\*
\end{longtable}

\endgroup

\subsection{Coverage accounting and interval uncertainty}

The registered study yields \RegisteredJointAvailable{} numerically
completed fits. Its \RegisteredJointFailures{} unavailable fits
remain in the coverage denominator. For a coordinate group $G$ in
setting $a$ and replication $r$, define
\[
 C_{ar,G}=\frac{1}{|G|}\sum_{k\in G}
  1\{\text{interval }I_{ark}\text{ is available}\}
  1\{\beta_k^\star\in I_{ark}\}.
\]
Group coverage is the mean of $C_{ar,G}$ over the registered
replications, with Monte Carlo standard error
$\operatorname{sd}(C_{ar,G})/\sqrt{B_a}$. This calculation retains
the dependence between coordinates in a dataset. When settings are
averaged within an input category, each setting receives equal weight
and its Monte Carlo variance contributes with the squared averaging
weight.

Table~\ref{tab:registered-joint-coverage} reports every setting's
active and null coverage. Table~\ref{tab:registered-joint-length}
gives mean lengths among available intervals; its Monte Carlo
standard error uses the joint replication averages of available
length and interval availability. At the individual-coordinate level,
Table~\ref{tab:registered-joint-coordinates} reports exact 95\%
binomial intervals for coverage. The standard error ratio divides
the mean reported standard error by the empirical standard deviation
among available coefficient estimates for that coordinate.

\begingroup\small

\begin{longtable}[t]{llll}
\caption{\label{tab:registered-joint-coverage}Registered full target coverage with Monte Carlo standard errors}\\
\toprule
Setting & Fits & Active & Null\\
\midrule
\endfirsthead
\caption[]{Registered full target coverage with Monte Carlo standard errors \textit{(continued)}}\\
\toprule
Setting & Fits & Active & Null\\
\midrule
\endhead

\endfoot
\bottomrule
\endlastfoot
A01 & 1000/1000 & 0.970 (0.004) & 0.972 (0.004)\\
A02 & 1000/1000 & 0.959 (0.004) & 0.974 (0.004)\\
A03 & 1000/1000 & 0.963 (0.004) & 0.968 (0.004)\\
A04 & 1000/1000 & 0.965 (0.004) & 0.975 (0.004)\\
A05 & 1000/1000 & 0.968 (0.004) & 0.977 (0.003)\\
A06 & 1000/1000 & 0.959 (0.005) & 0.974 (0.004)\\
A07 & 500/500 & 0.955 (0.007) & 0.973 (0.005)\\
A08 & 500/500 & 0.964 (0.006) & 0.979 (0.004)\\
A09 & 500/500 & 0.959 (0.006) & 0.971 (0.005)\\
A10 & 1000/1000 & 0.961 (0.004) & 0.964 (0.004)\\
A11 & 1000/1000 & 0.984 (0.003) & 0.983 (0.003)\\
A12 & 1000/1000 & 0.966 (0.004) & 0.980 (0.003)\\
B01 & 1000/1000 & 0.977 (0.003) & 0.980 (0.003)\\
B02 & 1000/1000 & 0.981 (0.003) & 0.984 (0.003)\\
B03 & 1000/1000 & 0.984 (0.003) & 0.988 (0.002)\\
B04 & 1000/1000 & 0.974 (0.004) & 0.983 (0.003)\\
B05 & 1000/1000 & 0.974 (0.004) & 0.982 (0.003)\\
B06 & 1000/1000 & 0.977 (0.003) & 0.981 (0.003)\\
B07 & 500/500 & 0.973 (0.005) & 0.985 (0.004)\\
B08 & 500/500 & 0.970 (0.005) & 0.983 (0.004)\\
B09 & 500/500 & 0.980 (0.004) & 0.982 (0.005)\\
B10 & 500/500 & 0.978 (0.005) & 0.985 (0.004)\\
B11 & 1000/1000 & 0.986 (0.003) & 0.988 (0.003)\\
B12 & 1000/1000 & 0.983 (0.003) & 0.988 (0.002)\\
B13 & 1000/1000 & 0.983 (0.003) & 0.988 (0.002)\\
B14 & 500/500 & 0.978 (0.005) & 0.976 (0.005)\\
B15 & 500/500 & 0.979 (0.004) & 0.984 (0.004)\\
B16 & 500/500 & 0.965 (0.006) & 0.960 (0.006)\\
B17 & 500/500 & 0.978 (0.005) & 0.984 (0.004)\\
B18 & 500/500 & 0.984 (0.004) & 0.972 (0.006)\\
B19 & 499/500 & 0.940 (0.008) & 0.940 (0.008)\\
C01 & 500/500 & 0.978 (0.005) & 0.977 (0.005)\\
C02 & 1000/1000 & 0.977 (0.003) & 0.978 (0.003)\\
C03 & 1000/1000 & 0.981 (0.003) & 0.976 (0.003)\\
C04 & 500/500 & 0.973 (0.005) & 0.976 (0.005)\\
C05 & 500/500 & 0.969 (0.005) & 0.980 (0.004)\\
C06 & 500/500 & 0.956 (0.006) & 0.967 (0.006)\\
C07 & 500/500 & 0.954 (0.007) & 0.954 (0.007)\\
C08 & 500/500 & 0.977 (0.005) & 0.984 (0.004)\\
C09 & 500/500 & 0.979 (0.004) & 0.981 (0.005)\\
C10 & 500/500 & 0.977 (0.005) & 0.983 (0.005)\\
C11 & 500/500 & 0.991 (0.003) & 0.984 (0.004)\\
C12 & 500/500 & 0.974 (0.005) & 0.984 (0.004)\\
C13 & 500/500 & 0.983 (0.004) & 0.980 (0.004)\\
C14 & 500/500 & 0.963 (0.006) & 0.957 (0.007)\\
C15 & 500/500 & 0.958 (0.007) & 0.945 (0.008)\\
C16 & 500/500 & 0.947 (0.008) & 0.961 (0.006)\\
D01 & 500/500 & 0.963 (0.006) & 0.974 (0.005)\\
D02 & 500/500 & 0.956 (0.007) & 0.975 (0.005)\\
D03 & 495/500 & 0.971 (0.006) & 0.972 (0.006)\\
D04 & 500/500 & 0.970 (0.005) & 0.977 (0.005)\\
D05 & 500/500 & 0.979 (0.005) & 0.992 (0.003)\\
D06 & 500/500 & 0.982 (0.004) & 0.983 (0.004)\\
E01 & 500/500 & 0.981 (0.005) & 0.982 (0.004)\\
E02 & 500/500 & 0.976 (0.005) & 0.977 (0.005)\\
E03 & 500/500 & 0.973 (0.005) & 0.982 (0.004)\\
E04 & 500/500 & 0.981 (0.005) & 0.973 (0.005)\\
E05 & 500/500 & 0.972 (0.005) & 0.986 (0.004)\\
E06 & 500/500 & 0.986 (0.004) & 0.987 (0.004)\\
E07 & 500/500 & 0.952 (0.007) & 0.963 (0.006)\\
E08 & 500/500 & 0.976 (0.005) & 0.993 (0.003)\\
F01 & 500/500 & 0.941 (0.007) & 0.965 (0.006)\\
F02 & 500/500 & 0.968 (0.006) & 0.987 (0.004)\\
F03 & 500/500 & 0.964 (0.006) & 0.974 (0.005)\\
F04 & 500/500 & 0.981 (0.004) & 0.979 (0.005)\\
F05 & 500/500 & 0.661 (0.012) & 0.694 (0.012)\\
F06 & 500/500 & 0.977 (0.005) & 0.989 (0.003)\\
H01 & 500/500 & 0.969 (0.005) & 0.975 (0.005)\\
H02 & 500/500 & 0.965 (0.006) & 0.974 (0.005)\\
H03 & 500/500 & 0.958 (0.007) & 0.966 (0.006)\\
H04 & 500/500 & 0.963 (0.006) & 0.972 (0.005)\\*
\end{longtable}


\begin{longtable}[t]{lll}
\caption{\label{tab:registered-joint-length}Registered interval lengths with Monte Carlo standard errors}\\
\toprule
Setting & Active & Null\\
\midrule
\endfirsthead
\caption[]{Registered interval lengths with Monte Carlo standard errors \textit{(continued)}}\\
\toprule
Setting & Active & Null\\
\midrule
\endhead

\endfoot
\bottomrule
\endlastfoot
A01 & 0.546 (0.006) & 0.632 (0.006)\\
A02 & 0.302 (0.002) & 0.343 (0.002)\\
A03 & 0.191 (0.001) & 0.212 (0.001)\\
A04 & 0.547 (0.005) & 0.641 (0.007)\\
A05 & 0.307 (0.002) & 0.352 (0.002)\\
A06 & 0.193 (0.001) & 0.215 (0.001)\\
A07 & 0.565 (0.008) & 0.657 (0.009)\\
A08 & 0.320 (0.003) & 0.369 (0.004)\\
A09 & 0.194 (0.001) & 0.217 (0.001)\\
A10 & 0.690 (0.008) & 0.823 (0.010)\\
A11 & 0.355 (0.003) & 0.409 (0.003)\\
A12 & 0.201 (0.001) & 0.226 (0.001)\\
B01 & 0.386 (0.003) & 0.438 (0.003)\\
B02 & 0.339 (0.002) & 0.381 (0.003)\\
B03 & 0.389 (0.003) & 0.445 (0.003)\\
B04 & 0.289 (0.002) & 0.331 (0.002)\\
B05 & 0.254 (0.002) & 0.289 (0.002)\\
B06 & 0.288 (0.002) & 0.332 (0.002)\\
B07 & 0.323 (0.003) & 0.379 (0.004)\\
B08 & 0.326 (0.003) & 0.371 (0.004)\\
B09 & 0.279 (0.003) & 0.315 (0.003)\\
B10 & 0.398 (0.004) & 0.450 (0.005)\\
B11 & 0.471 (0.004) & 0.537 (0.005)\\
B12 & 0.398 (0.003) & 0.452 (0.003)\\
B13 & 0.481 (0.004) & 0.538 (0.004)\\
B14 & 0.275 (0.003) & 0.314 (0.003)\\
B15 & 0.276 (0.003) & 0.320 (0.003)\\
B16 & 0.317 (0.004) & 0.335 (0.004)\\
B17 & 0.361 (0.004) & 0.401 (0.004)\\
B18 & 0.262 (0.002) & 0.297 (0.003)\\
B19 & 0.248 (0.003) & 0.265 (0.003)\\
C01 & 0.258 (0.002) & 0.292 (0.002)\\
C02 & 0.614 (0.005) & 0.696 (0.006)\\
C03 & 0.328 (0.002) & 0.370 (0.002)\\
C04 & 0.314 (0.003) & 0.355 (0.003)\\
C05 & 0.333 (0.003) & 0.374 (0.003)\\
C06 & 0.288 (0.003) & 0.308 (0.003)\\
C07 & 0.288 (0.003) & 0.307 (0.003)\\
C08 & 0.395 (0.004) & 0.438 (0.004)\\
C09 & 0.268 (0.002) & 0.300 (0.003)\\
C10 & 0.309 (0.003) & 0.352 (0.004)\\
C11 & 0.331 (0.004) & 0.367 (0.003)\\
C12 & 0.278 (0.003) & 0.320 (0.003)\\
C13 & 0.550 (0.006) & 0.627 (0.008)\\
C14 & 0.290 (0.003) & 0.298 (0.003)\\
C15 & 0.286 (0.003) & 0.293 (0.003)\\
C16 & 0.286 (0.003) & 0.307 (0.003)\\
D01 & 0.245 (0.003) & 0.249 (0.002)\\
D02 & 0.519 (0.005) & 0.617 (0.006)\\
D03 & 0.538 (0.005) & 0.544 (0.005)\\
D04 & 0.270 (0.002) & 0.288 (0.003)\\
D05 & 0.297 (0.003) & 0.318 (0.003)\\
D06 & 0.359 (0.004) & 0.360 (0.003)\\
E01 & 0.349 (0.004) & 0.408 (0.004)\\
E02 & 0.328 (0.003) & 0.365 (0.003)\\
E03 & 0.281 (0.002) & 0.321 (0.003)\\
E04 & 0.302 (0.002) & 0.338 (0.003)\\
E05 & 0.384 (0.004) & 0.449 (0.005)\\
E06 & 0.439 (0.006) & 0.501 (0.006)\\
E07 & 0.318 (0.003) & 0.359 (0.004)\\
E08 & 0.353 (0.004) & 0.422 (0.005)\\
F01 & 0.709 (0.011) & 0.775 (0.012)\\
F02 & 0.399 (0.004) & 0.423 (0.005)\\
F03 & 0.465 (0.005) & 0.488 (0.006)\\
F04 & 0.354 (0.003) & 0.363 (0.003)\\
F05 & 0.417 (0.005) & 0.449 (0.005)\\
F06 & 0.417 (0.004) & 0.477 (0.005)\\
H01 & 0.342 (0.003) & 0.387 (0.004)\\
H02 & 0.311 (0.003) & 0.357 (0.003)\\
H03 & 0.367 (0.004) & 0.407 (0.004)\\
H04 & 0.342 (0.003) & 0.395 (0.004)\\*
\end{longtable}


\begin{longtable}[t]{lrllll}
\caption{\label{tab:registered-joint-coordinates}Coordinate coverage with exact binomial intervals and standard error ratios}\\
\toprule
Setting & Coordinate & Available & Coverage & Interval & $\overline{s}/\mathrm{SD}$\\
\midrule
\endfirsthead
\caption[]{Coordinate coverage with exact binomial intervals and standard error ratios \textit{(continued)}}\\
\toprule
Setting & Coordinate & Available & Coverage & Interval & $\overline{s}/\mathrm{SD}$\\
\midrule
\endhead

\endfoot
\bottomrule
\endlastfoot
A01 & 1 & 1000/1000 & 0.964 (0.006) & {}[0.951, 0.975] & 1.439\\
A01 & 2 & 1000/1000 & 0.976 (0.005) & {}[0.964, 0.985] & 1.451\\
A01 & 6 & 1000/1000 & 0.974 (0.005) & {}[0.962, 0.983] & 1.542\\
A01 & 7 & 1000/1000 & 0.970 (0.005) & {}[0.957, 0.980] & 1.551\\
A02 & 1 & 1000/1000 & 0.954 (0.007) & {}[0.939, 0.966] & 1.202\\
A02 & 2 & 1000/1000 & 0.964 (0.006) & {}[0.951, 0.975] & 1.230\\
A02 & 6 & 1000/1000 & 0.968 (0.006) & {}[0.955, 0.978] & 1.256\\
A02 & 7 & 1000/1000 & 0.979 (0.005) & {}[0.968, 0.987] & 1.322\\
A03 & 1 & 1000/1000 & 0.964 (0.006) & {}[0.951, 0.975] & 1.109\\
A03 & 2 & 1000/1000 & 0.962 (0.006) & {}[0.948, 0.973] & 1.123\\
A03 & 6 & 1000/1000 & 0.967 (0.006) & {}[0.954, 0.977] & 1.155\\
A03 & 7 & 1000/1000 & 0.969 (0.005) & {}[0.956, 0.979] & 1.171\\
A04 & 1 & 1000/1000 & 0.963 (0.006) & {}[0.949, 0.974] & 1.400\\
A04 & 2 & 1000/1000 & 0.966 (0.006) & {}[0.953, 0.976] & 1.400\\
A04 & 6 & 1000/1000 & 0.974 (0.005) & {}[0.962, 0.983] & 1.423\\
A04 & 7 & 1000/1000 & 0.975 (0.005) & {}[0.963, 0.984] & 1.511\\
A05 & 1 & 1000/1000 & 0.965 (0.006) & {}[0.952, 0.976] & 1.222\\
A05 & 2 & 1000/1000 & 0.971 (0.005) & {}[0.959, 0.980] & 1.290\\
A05 & 6 & 1000/1000 & 0.974 (0.005) & {}[0.962, 0.983] & 1.359\\
A05 & 7 & 1000/1000 & 0.979 (0.005) & {}[0.968, 0.987] & 1.361\\
A06 & 1 & 1000/1000 & 0.959 (0.006) & {}[0.945, 0.970] & 1.095\\
A06 & 2 & 1000/1000 & 0.959 (0.006) & {}[0.945, 0.970] & 1.152\\
A06 & 6 & 1000/1000 & 0.974 (0.005) & {}[0.962, 0.983] & 1.178\\
A06 & 7 & 1000/1000 & 0.974 (0.005) & {}[0.962, 0.983] & 1.231\\
A07 & 1 & 500/500 & 0.950 (0.010) & {}[0.927, 0.967] & 1.398\\
A07 & 2 & 500/500 & 0.960 (0.009) & {}[0.939, 0.975] & 1.463\\
A07 & 6 & 500/500 & 0.974 (0.007) & {}[0.956, 0.986] & 1.508\\
A07 & 7 & 500/500 & 0.972 (0.007) & {}[0.953, 0.985] & 1.524\\
A08 & 1 & 500/500 & 0.968 (0.008) & {}[0.949, 0.982] & 1.345\\
A08 & 2 & 500/500 & 0.960 (0.009) & {}[0.939, 0.975] & 1.288\\
A08 & 6 & 500/500 & 0.984 (0.006) & {}[0.969, 0.993] & 1.459\\
A08 & 7 & 500/500 & 0.974 (0.007) & {}[0.956, 0.986] & 1.481\\
A09 & 1 & 500/500 & 0.964 (0.008) & {}[0.944, 0.979] & 1.132\\
A09 & 2 & 500/500 & 0.954 (0.009) & {}[0.932, 0.971] & 1.140\\
A09 & 6 & 500/500 & 0.966 (0.008) & {}[0.946, 0.980] & 1.177\\
A09 & 7 & 500/500 & 0.976 (0.007) & {}[0.958, 0.988] & 1.234\\
A10 & 1 & 1000/1000 & 0.960 (0.006) & {}[0.946, 0.971] & 1.439\\
A10 & 2 & 1000/1000 & 0.962 (0.006) & {}[0.948, 0.973] & 1.386\\
A10 & 11 & 1000/1000 & 0.965 (0.006) & {}[0.952, 0.976] & 1.449\\
A10 & 12 & 1000/1000 & 0.963 (0.006) & {}[0.949, 0.974] & 1.561\\
A11 & 1 & 1000/1000 & 0.988 (0.003) & {}[0.979, 0.994] & 1.401\\
A11 & 2 & 1000/1000 & 0.979 (0.005) & {}[0.968, 0.987] & 1.413\\
A11 & 11 & 1000/1000 & 0.983 (0.004) & {}[0.973, 0.990] & 1.492\\
A11 & 12 & 1000/1000 & 0.983 (0.004) & {}[0.973, 0.990] & 1.519\\
A12 & 1 & 1000/1000 & 0.964 (0.006) & {}[0.951, 0.975] & 1.174\\
A12 & 2 & 1000/1000 & 0.967 (0.006) & {}[0.954, 0.977] & 1.193\\
A12 & 11 & 1000/1000 & 0.979 (0.005) & {}[0.968, 0.987] & 1.237\\
A12 & 12 & 1000/1000 & 0.980 (0.004) & {}[0.969, 0.988] & 1.225\\
B01 & 1 & 1000/1000 & 0.975 (0.005) & {}[0.963, 0.984] & 1.379\\
B01 & 2 & 1000/1000 & 0.978 (0.005) & {}[0.967, 0.986] & 1.403\\
B01 & 9 & 1000/1000 & 0.981 (0.004) & {}[0.970, 0.989] & 1.486\\
B01 & 10 & 1000/1000 & 0.979 (0.005) & {}[0.968, 0.987] & 1.493\\
B02 & 1 & 1000/1000 & 0.981 (0.004) & {}[0.970, 0.989] & 1.391\\
B02 & 2 & 1000/1000 & 0.980 (0.004) & {}[0.969, 0.988] & 1.384\\
B02 & 9 & 1000/1000 & 0.982 (0.004) & {}[0.972, 0.989] & 1.419\\
B02 & 10 & 1000/1000 & 0.985 (0.004) & {}[0.975, 0.992] & 1.457\\
B03 & 1 & 1000/1000 & 0.980 (0.004) & {}[0.969, 0.988] & 1.490\\
B03 & 2 & 1000/1000 & 0.988 (0.003) & {}[0.979, 0.994] & 1.496\\
B03 & 9 & 1000/1000 & 0.989 (0.003) & {}[0.980, 0.994] & 1.598\\
B03 & 10 & 1000/1000 & 0.987 (0.004) & {}[0.978, 0.993] & 1.549\\
B04 & 1 & 1000/1000 & 0.973 (0.005) & {}[0.961, 0.982] & 1.307\\
B04 & 2 & 1000/1000 & 0.975 (0.005) & {}[0.963, 0.984] & 1.332\\
B04 & 9 & 1000/1000 & 0.977 (0.005) & {}[0.966, 0.985] & 1.381\\
B04 & 10 & 1000/1000 & 0.989 (0.003) & {}[0.980, 0.994] & 1.448\\
B05 & 1 & 1000/1000 & 0.966 (0.006) & {}[0.953, 0.976] & 1.277\\
B05 & 2 & 1000/1000 & 0.982 (0.004) & {}[0.972, 0.989] & 1.325\\
B05 & 9 & 1000/1000 & 0.980 (0.004) & {}[0.969, 0.988] & 1.345\\
B05 & 10 & 1000/1000 & 0.984 (0.004) & {}[0.974, 0.991] & 1.396\\
B06 & 1 & 1000/1000 & 0.977 (0.005) & {}[0.966, 0.985] & 1.309\\
B06 & 2 & 1000/1000 & 0.977 (0.005) & {}[0.966, 0.985] & 1.340\\
B06 & 9 & 1000/1000 & 0.979 (0.005) & {}[0.968, 0.987] & 1.395\\
B06 & 10 & 1000/1000 & 0.982 (0.004) & {}[0.972, 0.989] & 1.428\\
B07 & 1 & 500/500 & 0.980 (0.006) & {}[0.964, 0.990] & 1.401\\
B07 & 2 & 500/500 & 0.966 (0.008) & {}[0.946, 0.980] & 1.292\\
B07 & 9 & 500/500 & 0.978 (0.007) & {}[0.961, 0.989] & 1.457\\
B07 & 10 & 500/500 & 0.992 (0.004) & {}[0.980, 0.998] & 1.528\\
B08 & 1 & 500/500 & 0.970 (0.008) & {}[0.951, 0.983] & 1.331\\
B08 & 2 & 500/500 & 0.970 (0.008) & {}[0.951, 0.983] & 1.385\\
B08 & 9 & 500/500 & 0.980 (0.006) & {}[0.964, 0.990] & 1.518\\
B08 & 10 & 500/500 & 0.986 (0.005) & {}[0.971, 0.994] & 1.489\\
B09 & 1 & 500/500 & 0.974 (0.007) & {}[0.956, 0.986] & 1.396\\
B09 & 2 & 500/500 & 0.986 (0.005) & {}[0.971, 0.994] & 1.276\\
B09 & 9 & 500/500 & 0.980 (0.006) & {}[0.964, 0.990] & 1.379\\
B09 & 10 & 500/500 & 0.984 (0.006) & {}[0.969, 0.993] & 1.394\\
B10 & 1 & 500/500 & 0.980 (0.006) & {}[0.964, 0.990] & 1.468\\
B10 & 2 & 500/500 & 0.976 (0.007) & {}[0.958, 0.988] & 1.471\\
B10 & 9 & 500/500 & 0.988 (0.005) & {}[0.974, 0.996] & 1.487\\
B10 & 10 & 500/500 & 0.982 (0.006) & {}[0.966, 0.992] & 1.599\\
B11 & 1 & 1000/1000 & 0.987 (0.004) & {}[0.978, 0.993] & 1.529\\
B11 & 2 & 1000/1000 & 0.984 (0.004) & {}[0.974, 0.991] & 1.526\\
B11 & 9 & 1000/1000 & 0.989 (0.003) & {}[0.980, 0.994] & 1.658\\
B11 & 10 & 1000/1000 & 0.987 (0.004) & {}[0.978, 0.993] & 1.619\\
B12 & 1 & 1000/1000 & 0.980 (0.004) & {}[0.969, 0.988] & 1.377\\
B12 & 2 & 1000/1000 & 0.985 (0.004) & {}[0.975, 0.992] & 1.483\\
B12 & 9 & 1000/1000 & 0.987 (0.004) & {}[0.978, 0.993] & 1.550\\
B12 & 10 & 1000/1000 & 0.988 (0.003) & {}[0.979, 0.994] & 1.547\\
B13 & 1 & 1000/1000 & 0.981 (0.004) & {}[0.970, 0.989] & 1.521\\
B13 & 2 & 1000/1000 & 0.984 (0.004) & {}[0.974, 0.991] & 1.489\\
B13 & 9 & 1000/1000 & 0.987 (0.004) & {}[0.978, 0.993] & 1.583\\
B13 & 10 & 1000/1000 & 0.989 (0.003) & {}[0.980, 0.994] & 1.610\\
B14 & 1 & 500/500 & 0.976 (0.007) & {}[0.958, 0.988] & 1.305\\
B14 & 2 & 500/500 & 0.980 (0.006) & {}[0.964, 0.990] & 1.332\\
B14 & 9 & 500/500 & 0.972 (0.007) & {}[0.953, 0.985] & 1.368\\
B14 & 10 & 500/500 & 0.980 (0.006) & {}[0.964, 0.990] & 1.427\\
B15 & 1 & 500/500 & 0.980 (0.006) & {}[0.964, 0.990] & 1.244\\
B15 & 2 & 500/500 & 0.978 (0.007) & {}[0.961, 0.989] & 1.334\\
B15 & 9 & 500/500 & 0.986 (0.005) & {}[0.971, 0.994] & 1.476\\
B15 & 10 & 500/500 & 0.982 (0.006) & {}[0.966, 0.992] & 1.423\\
B16 & 1 & 500/500 & 0.964 (0.008) & {}[0.944, 0.979] & 1.145\\
B16 & 2 & 500/500 & 0.966 (0.008) & {}[0.946, 0.980] & 1.204\\
B16 & 9 & 500/500 & 0.954 (0.009) & {}[0.932, 0.971] & 1.175\\
B16 & 10 & 500/500 & 0.966 (0.008) & {}[0.946, 0.980] & 1.210\\
B17 & 1 & 500/500 & 0.972 (0.007) & {}[0.953, 0.985] & 1.416\\
B17 & 2 & 500/500 & 0.984 (0.006) & {}[0.969, 0.993] & 1.427\\
B17 & 9 & 500/500 & 0.980 (0.006) & {}[0.964, 0.990] & 1.453\\
B17 & 10 & 500/500 & 0.988 (0.005) & {}[0.974, 0.996] & 1.599\\
B18 & 1 & 500/500 & 0.988 (0.005) & {}[0.974, 0.996] & 1.358\\
B18 & 2 & 500/500 & 0.980 (0.006) & {}[0.964, 0.990] & 1.314\\
B18 & 9 & 500/500 & 0.964 (0.008) & {}[0.944, 0.979] & 1.313\\
B18 & 10 & 500/500 & 0.980 (0.006) & {}[0.964, 0.990] & 1.344\\
B19 & 1 & 499/500 & 0.948 (0.010) & {}[0.925, 0.966] & 1.102\\
B19 & 2 & 499/500 & 0.932 (0.011) & {}[0.906, 0.952] & 1.099\\
B19 & 9 & 499/500 & 0.936 (0.011) & {}[0.911, 0.956] & 1.172\\
B19 & 10 & 499/500 & 0.944 (0.010) & {}[0.920, 0.962] & 1.131\\
C01 & 1 & 500/500 & 0.966 (0.008) & {}[0.946, 0.980] & 1.163\\
C01 & 2 & 500/500 & 0.990 (0.004) & {}[0.977, 0.997] & 1.293\\
C01 & 9 & 500/500 & 0.968 (0.008) & {}[0.949, 0.982] & 1.257\\
C01 & 10 & 500/500 & 0.986 (0.005) & {}[0.971, 0.994] & 1.349\\
C02 & 1 & 1000/1000 & 0.974 (0.005) & {}[0.962, 0.983] & 1.455\\
C02 & 2 & 1000/1000 & 0.980 (0.004) & {}[0.969, 0.988] & 1.522\\
C02 & 9 & 1000/1000 & 0.981 (0.004) & {}[0.970, 0.989] & 1.517\\
C02 & 10 & 1000/1000 & 0.974 (0.005) & {}[0.962, 0.983] & 1.486\\
C03 & 1 & 1000/1000 & 0.980 (0.004) & {}[0.969, 0.988] & 1.305\\
C03 & 2 & 1000/1000 & 0.982 (0.004) & {}[0.972, 0.989] & 1.350\\
C03 & 9 & 1000/1000 & 0.975 (0.005) & {}[0.963, 0.984] & 1.340\\
C03 & 10 & 1000/1000 & 0.977 (0.005) & {}[0.966, 0.985] & 1.378\\
C04 & 1 & 500/500 & 0.970 (0.008) & {}[0.951, 0.983] & 1.258\\
C04 & 2 & 500/500 & 0.976 (0.007) & {}[0.958, 0.988] & 1.280\\
C04 & 9 & 500/500 & 0.972 (0.007) & {}[0.953, 0.985] & 1.305\\
C04 & 10 & 500/500 & 0.980 (0.006) & {}[0.964, 0.990] & 1.439\\
C05 & 1 & 500/500 & 0.970 (0.008) & {}[0.951, 0.983] & 1.259\\
C05 & 2 & 500/500 & 0.968 (0.008) & {}[0.949, 0.982] & 1.307\\
C05 & 9 & 500/500 & 0.978 (0.007) & {}[0.961, 0.989] & 1.468\\
C05 & 10 & 500/500 & 0.982 (0.006) & {}[0.966, 0.992] & 1.451\\
C06 & 1 & 500/500 & 0.956 (0.009) & {}[0.934, 0.972] & 1.191\\
C06 & 2 & 500/500 & 0.956 (0.009) & {}[0.934, 0.972] & 1.121\\
C06 & 9 & 500/500 & 0.968 (0.008) & {}[0.949, 0.982] & 1.241\\
C06 & 10 & 500/500 & 0.966 (0.008) & {}[0.946, 0.980] & 1.167\\
C07 & 1 & 500/500 & 0.946 (0.010) & {}[0.922, 0.964] & 1.158\\
C07 & 2 & 500/500 & 0.962 (0.009) & {}[0.941, 0.977] & 1.242\\
C07 & 9 & 500/500 & 0.952 (0.010) & {}[0.929, 0.969] & 1.210\\
C07 & 10 & 500/500 & 0.956 (0.009) & {}[0.934, 0.972] & 1.232\\
C08 & 1 & 500/500 & 0.974 (0.007) & {}[0.956, 0.986] & 1.407\\
C08 & 2 & 500/500 & 0.980 (0.006) & {}[0.964, 0.990] & 1.433\\
C08 & 9 & 500/500 & 0.990 (0.004) & {}[0.977, 0.997] & 1.594\\
C08 & 10 & 500/500 & 0.978 (0.007) & {}[0.961, 0.989] & 1.445\\
C09 & 1 & 500/500 & 0.984 (0.006) & {}[0.969, 0.993] & 1.262\\
C09 & 2 & 500/500 & 0.974 (0.007) & {}[0.956, 0.986] & 1.250\\
C09 & 9 & 500/500 & 0.986 (0.005) & {}[0.971, 0.994] & 1.352\\
C09 & 10 & 500/500 & 0.976 (0.007) & {}[0.958, 0.988] & 1.302\\
C10 & 1 & 500/500 & 0.978 (0.007) & {}[0.961, 0.989] & 1.268\\
C10 & 2 & 500/500 & 0.976 (0.007) & {}[0.958, 0.988] & 1.362\\
C10 & 9 & 500/500 & 0.976 (0.007) & {}[0.958, 0.988] & 1.351\\
C10 & 10 & 500/500 & 0.990 (0.004) & {}[0.977, 0.997] & 1.454\\
C11 & 1 & 500/500 & 0.988 (0.005) & {}[0.974, 0.996] & 1.432\\
C11 & 2 & 500/500 & 0.994 (0.003) & {}[0.983, 0.999] & 1.471\\
C11 & 9 & 500/500 & 0.980 (0.006) & {}[0.964, 0.990] & 1.335\\
C11 & 10 & 500/500 & 0.988 (0.005) & {}[0.974, 0.996] & 1.443\\
C12 & 1 & 500/500 & 0.984 (0.006) & {}[0.969, 0.993] & 1.369\\
C12 & 2 & 500/500 & 0.964 (0.008) & {}[0.944, 0.979] & 1.325\\
C12 & 9 & 500/500 & 0.982 (0.006) & {}[0.966, 0.992] & 1.388\\
C12 & 10 & 500/500 & 0.986 (0.005) & {}[0.971, 0.994] & 1.381\\
C13 & 1 & 500/500 & 0.984 (0.006) & {}[0.969, 0.993] & 1.526\\
C13 & 2 & 500/500 & 0.982 (0.006) & {}[0.966, 0.992] & 1.557\\
C13 & 9 & 500/500 & 0.986 (0.005) & {}[0.971, 0.994] & 1.538\\
C13 & 10 & 500/500 & 0.974 (0.007) & {}[0.956, 0.986] & 1.543\\
C14 & 1 & 500/500 & 0.954 (0.009) & {}[0.932, 0.971] & 1.112\\
C14 & 2 & 500/500 & 0.972 (0.007) & {}[0.953, 0.985] & 1.255\\
C14 & 9 & 500/500 & 0.954 (0.009) & {}[0.932, 0.971] & 1.130\\
C14 & 10 & 500/500 & 0.960 (0.009) & {}[0.939, 0.975] & 1.198\\
C15 & 1 & 500/500 & 0.958 (0.009) & {}[0.937, 0.974] & 1.159\\
C15 & 2 & 500/500 & 0.958 (0.009) & {}[0.937, 0.974] & 1.193\\
C15 & 9 & 500/500 & 0.938 (0.011) & {}[0.913, 0.957] & 1.100\\
C15 & 10 & 500/500 & 0.952 (0.010) & {}[0.929, 0.969] & 1.145\\
C16 & 1 & 500/500 & 0.946 (0.010) & {}[0.922, 0.964] & 1.102\\
C16 & 2 & 500/500 & 0.948 (0.010) & {}[0.925, 0.966] & 1.115\\
C16 & 9 & 500/500 & 0.964 (0.008) & {}[0.944, 0.979] & 1.252\\
C16 & 10 & 500/500 & 0.958 (0.009) & {}[0.937, 0.974] & 1.150\\
D01 & 1 & 500/500 & 0.968 (0.008) & {}[0.949, 0.982] & 1.272\\
D01 & 2 & 500/500 & 0.958 (0.009) & {}[0.937, 0.974] & 1.171\\
D01 & 9 & 500/500 & 0.978 (0.007) & {}[0.961, 0.989] & 1.273\\
D01 & 10 & 500/500 & 0.970 (0.008) & {}[0.951, 0.983] & 1.335\\
D02 & 1 & 500/500 & 0.954 (0.009) & {}[0.932, 0.971] & 1.222\\
D02 & 2 & 500/500 & 0.958 (0.009) & {}[0.937, 0.974] & 1.210\\
D02 & 9 & 500/500 & 0.972 (0.007) & {}[0.953, 0.985] & 1.286\\
D02 & 10 & 500/500 & 0.978 (0.007) & {}[0.961, 0.989] & 1.361\\
D03 & 1 & 495/500 & 0.968 (0.008) & {}[0.949, 0.982] & 1.473\\
D03 & 2 & 495/500 & 0.974 (0.007) & {}[0.956, 0.986] & 1.501\\
D03 & 9 & 495/500 & 0.972 (0.007) & {}[0.953, 0.985] & 1.523\\
D03 & 10 & 495/500 & 0.972 (0.007) & {}[0.953, 0.985] & 1.401\\
D04 & 1 & 500/500 & 0.968 (0.008) & {}[0.949, 0.982] & 1.303\\
D04 & 2 & 500/500 & 0.972 (0.007) & {}[0.953, 0.985] & 1.254\\
D04 & 9 & 500/500 & 0.982 (0.006) & {}[0.966, 0.992] & 1.279\\
D04 & 10 & 500/500 & 0.972 (0.007) & {}[0.953, 0.985] & 1.284\\
D05 & 1 & 500/500 & 0.978 (0.007) & {}[0.961, 0.989] & 1.446\\
D05 & 2 & 500/500 & 0.980 (0.006) & {}[0.964, 0.990] & 1.413\\
D05 & 9 & 500/500 & 0.992 (0.004) & {}[0.980, 0.998] & 1.535\\
D05 & 10 & 500/500 & 0.992 (0.004) & {}[0.980, 0.998] & 1.544\\
D06 & 1 & 500/500 & 0.988 (0.005) & {}[0.974, 0.996] & 1.476\\
D06 & 2 & 500/500 & 0.976 (0.007) & {}[0.958, 0.988] & 1.364\\
D06 & 9 & 500/500 & 0.980 (0.006) & {}[0.964, 0.990] & 1.466\\
D06 & 10 & 500/500 & 0.986 (0.005) & {}[0.971, 0.994] & 1.377\\
E01 & 1 & 500/500 & 0.976 (0.007) & {}[0.958, 0.988] & 1.462\\
E01 & 2 & 500/500 & 0.986 (0.005) & {}[0.971, 0.994] & 1.486\\
E01 & 9 & 500/500 & 0.982 (0.006) & {}[0.966, 0.992] & 1.470\\
E01 & 10 & 500/500 & 0.982 (0.006) & {}[0.966, 0.992] & 1.603\\
E02 & 1 & 500/500 & 0.972 (0.007) & {}[0.953, 0.985] & 1.298\\
E02 & 2 & 500/500 & 0.980 (0.006) & {}[0.964, 0.990] & 1.290\\
E02 & 9 & 500/500 & 0.980 (0.006) & {}[0.964, 0.990] & 1.336\\
E02 & 10 & 500/500 & 0.974 (0.007) & {}[0.956, 0.986] & 1.322\\
E03 & 1 & 500/500 & 0.980 (0.006) & {}[0.964, 0.990] & 1.246\\
E03 & 2 & 500/500 & 0.966 (0.008) & {}[0.946, 0.980] & 1.284\\
E03 & 9 & 500/500 & 0.980 (0.006) & {}[0.964, 0.990] & 1.362\\
E03 & 10 & 500/500 & 0.984 (0.006) & {}[0.969, 0.993] & 1.411\\
E04 & 1 & 500/500 & 0.978 (0.007) & {}[0.961, 0.989] & 1.330\\
E04 & 2 & 500/500 & 0.984 (0.006) & {}[0.969, 0.993] & 1.351\\
E04 & 9 & 500/500 & 0.970 (0.008) & {}[0.951, 0.983] & 1.312\\
E04 & 10 & 500/500 & 0.976 (0.007) & {}[0.958, 0.988] & 1.439\\
E05 & 1 & 500/500 & 0.964 (0.008) & {}[0.944, 0.979] & 1.401\\
E05 & 2 & 500/500 & 0.980 (0.006) & {}[0.964, 0.990] & 1.454\\
E05 & 9 & 500/500 & 0.986 (0.005) & {}[0.971, 0.994] & 1.510\\
E05 & 10 & 500/500 & 0.986 (0.005) & {}[0.971, 0.994] & 1.645\\
E06 & 1 & 500/500 & 0.988 (0.005) & {}[0.974, 0.996] & 1.666\\
E06 & 2 & 500/500 & 0.984 (0.006) & {}[0.969, 0.993] & 1.586\\
E06 & 9 & 500/500 & 0.982 (0.006) & {}[0.966, 0.992] & 1.664\\
E06 & 10 & 500/500 & 0.992 (0.004) & {}[0.980, 0.998] & 1.692\\
E07 & 1 & 500/500 & 0.954 (0.009) & {}[0.932, 0.971] & 1.220\\
E07 & 2 & 500/500 & 0.950 (0.010) & {}[0.927, 0.967] & 1.186\\
E07 & 9 & 500/500 & 0.960 (0.009) & {}[0.939, 0.975] & 1.304\\
E07 & 10 & 500/500 & 0.966 (0.008) & {}[0.946, 0.980] & 1.321\\
E08 & 1 & 500/500 & 0.976 (0.007) & {}[0.958, 0.988] & 1.397\\
E08 & 2 & 500/500 & 0.976 (0.007) & {}[0.958, 0.988] & 1.408\\
E08 & 9 & 500/500 & 0.992 (0.004) & {}[0.980, 0.998] & 1.574\\
E08 & 10 & 500/500 & 0.994 (0.003) & {}[0.983, 0.999] & 1.753\\
F01 & 1 & 500/500 & 0.966 (0.008) & {}[0.946, 0.980] & 1.486\\
F01 & 2 & 500/500 & 0.916 (0.012) & {}[0.888, 0.939] & 1.221\\
F01 & 11 & 500/500 & 0.968 (0.008) & {}[0.949, 0.982] & 1.526\\
F01 & 12 & 500/500 & 0.962 (0.009) & {}[0.941, 0.977] & 1.591\\
F02 & 1 & 500/500 & 0.984 (0.006) & {}[0.969, 0.993] & 1.368\\
F02 & 2 & 500/500 & 0.952 (0.010) & {}[0.929, 0.969] & 1.187\\
F02 & 11 & 500/500 & 0.990 (0.004) & {}[0.977, 0.997] & 1.480\\
F02 & 12 & 500/500 & 0.984 (0.006) & {}[0.969, 0.993] & 1.487\\
F03 & 1 & 500/500 & 0.972 (0.007) & {}[0.953, 0.985] & 1.354\\
F03 & 2 & 500/500 & 0.956 (0.009) & {}[0.934, 0.972] & 1.245\\
F03 & 11 & 500/500 & 0.976 (0.007) & {}[0.958, 0.988] & 1.406\\
F03 & 12 & 500/500 & 0.972 (0.007) & {}[0.953, 0.985] & 1.411\\
F04 & 1 & 500/500 & 0.980 (0.006) & {}[0.964, 0.990] & 1.424\\
F04 & 2 & 500/500 & 0.982 (0.006) & {}[0.966, 0.992] & 1.460\\
F04 & 11 & 500/500 & 0.980 (0.006) & {}[0.964, 0.990] & 1.417\\
F04 & 12 & 500/500 & 0.978 (0.007) & {}[0.961, 0.989] & 1.391\\
F05 & 1 & 500/500 & 0.958 (0.009) & {}[0.937, 0.974] & 1.290\\
F05 & 2 & 500/500 & 0.364 (0.022) & {}[0.322, 0.408] & 0.978\\
F05 & 11 & 500/500 & 0.424 (0.022) & {}[0.380, 0.469] & 1.142\\
F05 & 12 & 500/500 & 0.964 (0.008) & {}[0.944, 0.979] & 1.390\\
F06 & 1 & 500/500 & 0.986 (0.005) & {}[0.971, 0.994] & 1.410\\
F06 & 2 & 500/500 & 0.968 (0.008) & {}[0.949, 0.982] & 1.439\\
F06 & 11 & 500/500 & 0.988 (0.005) & {}[0.974, 0.996] & 1.544\\
F06 & 12 & 500/500 & 0.990 (0.004) & {}[0.977, 0.997] & 1.663\\
H01 & 1 & 500/500 & 0.972 (0.007) & {}[0.953, 0.985] & 1.247\\
H01 & 2 & 500/500 & 0.966 (0.008) & {}[0.946, 0.980] & 1.240\\
H01 & 9 & 500/500 & 0.972 (0.007) & {}[0.953, 0.985] & 1.310\\
H01 & 10 & 500/500 & 0.978 (0.007) & {}[0.961, 0.989] & 1.377\\
H02 & 1 & 500/500 & 0.968 (0.008) & {}[0.949, 0.982] & 1.230\\
H02 & 2 & 500/500 & 0.962 (0.009) & {}[0.941, 0.977] & 1.237\\
H02 & 9 & 500/500 & 0.970 (0.008) & {}[0.951, 0.983] & 1.331\\
H02 & 10 & 500/500 & 0.978 (0.007) & {}[0.961, 0.989] & 1.470\\
H03 & 1 & 500/500 & 0.948 (0.010) & {}[0.925, 0.966] & 1.209\\
H03 & 2 & 500/500 & 0.968 (0.008) & {}[0.949, 0.982] & 1.258\\
H03 & 9 & 500/500 & 0.968 (0.008) & {}[0.949, 0.982] & 1.337\\
H03 & 10 & 500/500 & 0.964 (0.008) & {}[0.944, 0.979] & 1.282\\
H04 & 1 & 500/500 & 0.954 (0.009) & {}[0.932, 0.971] & 1.200\\
H04 & 2 & 500/500 & 0.972 (0.007) & {}[0.953, 0.985] & 1.261\\
H04 & 9 & 500/500 & 0.976 (0.007) & {}[0.958, 0.988] & 1.472\\
H04 & 10 & 500/500 & 0.968 (0.008) & {}[0.949, 0.982] & 1.367\\*
\end{longtable}

\endgroup

\subsection{Target changes induced by variance filtering}

The Gaussian projection formula gives the selected population target
from the saved supports in F05 and F06. For each estimation fold,
we evaluate $\Sigma_{AA}^{-1}\Sigma_{A,S_0}\beta^0_{S_0}$ in the
original generating coordinate system and then combine the requested
coordinate with the estimator's fold weights.
Table~\ref{tab:registered-prescreen-targets} reports how often the
result differs from the full target by more than $10^{-8}$, with
exact binomial intervals and Monte Carlo errors.

Variance filtering changes the second-coordinate target in
\VarianceFilterSecondShifted{} of 500 datasets and the
eleventh-coordinate target in \VarianceFilterEleventhShifted{}.
Their conditional adjustment variables can be absent from the
retained pool. In F06, every requested selected target agrees with
its full-design coordinate within the stated tolerance. Both settings'
coverage in Table~\ref{tab:registered-joint-coverage} is evaluated
against the full target.

\begin{table}

\caption{\label{tab:registered-prescreen-targets}Gaussian target displacements after variance filtering and full design fitting}
\centering
\begin{tabular}[t]{lrllll}
\toprule
Setting & Coordinate & Shifted & Fraction & Interval & Displacement\\
\midrule
F05 & 1 & 0/500 & 0.000 (0.000) & {}[0.000, 0.007] & 0.000 (0.000)\\
F05 & 2 & 463/500 & 0.926 (0.012) & {}[0.899, 0.947] & 0.233 (0.003)\\
F05 & 11 & 472/500 & 0.944 (0.010) & {}[0.920, 0.962] & -0.237 (0.003)\\
F05 & 12 & 0/500 & 0.000 (0.000) & {}[0.000, 0.007] & 0.000 (0.000)\\
F06 & 1 & 0/500 & 0.000 (0.000) & {}[0.000, 0.007] & 0.000 (0.000)\\
F06 & 2 & 0/500 & 0.000 (0.000) & {}[0.000, 0.007] & 0.000 (0.000)\\
F06 & 11 & 0/500 & 0.000 (0.000) & {}[0.000, 0.007] & 0.000 (0.000)\\
F06 & 12 & 0/500 & 0.000 (0.000) & {}[0.000, 0.007] & 0.000 (0.000)\\
\bottomrule
\end{tabular}
\end{table}


\subsection{Generating distributions and independent replication}

Table~\ref{tab:joint-registry} specifies every setting in the independent
evaluation. In each setting, 60 datasets are generated with cluster
sizes sampled uniformly from $\{3,\ldots,8\}$. Rows of $X$ are
independent Gaussian with AR(1) covariance. The first five
coefficients are nonzero and alternate in sign. For a random
intercept, $b_i\sim N(0,1)$. With a time slope, the random-effects
covariance has diagonal $(1,0.25)$ and correlation $0.4$.
Innovations are independent Gaussian or $t_3/\sqrt3$, shifted by
their distribution's $\tau$th quantile. Covariates are independent
of the entire cluster nuisance and innovation vectors.

The random seed is $1800000000+10000a+r$, where $a$ is the
one-based configuration row and $r$ is the replication number.
The full method and paired ablation use the identical generated
sample and deterministic cluster partitions. The ablation removes
only the outcome-free conditional covariate support. Both use
$\Lambda=I$, $h=n^{-3/10}$ and normal 95\% intervals from their
paired cluster contributions. Two active and two null coordinates
are requested in every dataset.

\begin{longtable}{lrrrrrrr}
\caption{Simulation configurations with sixty independent datasets per row}\label{tab:joint-registry}\\
\toprule
Setting & $n$ & $p$ & $q$ & $\tau$ & Error & Signal & $\rho_X$ \\
\midrule\endfirsthead
\toprule
Setting & $n$ & $p$ & $q$ & $\tau$ & Error & Signal & $\rho_X$ \\
\midrule\endhead
N200\_Q1\_T25 & 200 &  500 & 1 & 0.25 & gaussian & 0.75 & 0.5 \\
N200\_Q1\_T50 & 200 &  500 & 1 & 0.50 & gaussian & 0.75 & 0.5 \\
N200\_Q1\_T75 & 200 &  500 & 1 & 0.75 & gaussian & 0.75 & 0.5 \\
N200\_Q2\_T25 & 200 &  500 & 2 & 0.25 & gaussian & 0.75 & 0.5 \\
N200\_Q2\_T50 & 200 &  500 & 2 & 0.50 & gaussian & 0.75 & 0.5 \\
N200\_Q2\_T75 & 200 &  500 & 2 & 0.75 & gaussian & 0.75 & 0.5 \\
N400\_Q1\_T25 & 400 &  500 & 1 & 0.25 & gaussian & 0.75 & 0.5 \\
N400\_Q1\_T50 & 400 &  500 & 1 & 0.50 & gaussian & 0.75 & 0.5 \\
N400\_Q1\_T75 & 400 &  500 & 1 & 0.75 & gaussian & 0.75 & 0.5 \\
N400\_Q2\_T25 & 400 &  500 & 2 & 0.25 & gaussian & 0.75 & 0.5 \\
N400\_Q2\_T50 & 400 &  500 & 2 & 0.50 & gaussian & 0.75 & 0.5 \\
N400\_Q2\_T75 & 400 &  500 & 2 & 0.75 & gaussian & 0.75 & 0.5 \\
T3\_Q1\_T25 & 200 &  500 & 1 & 0.25 & t3 & 0.75 & 0.5 \\
T3\_Q1\_T75 & 200 &  500 & 1 & 0.75 & t3 & 0.75 & 0.5 \\
T3\_Q2\_T25 & 200 &  500 & 2 & 0.25 & t3 & 0.75 & 0.5 \\
T3\_Q2\_T75 & 200 &  500 & 2 & 0.75 & t3 & 0.75 & 0.5 \\
WEAK\_Q1\_T25 & 200 &  500 & 1 & 0.25 & gaussian & 0.40 & 0.5 \\
WEAK\_Q1\_T50 & 200 &  500 & 1 & 0.50 & gaussian & 0.40 & 0.5 \\
WEAK\_Q2\_T25 & 200 &  500 & 2 & 0.25 & gaussian & 0.40 & 0.5 \\
CORR\_Q2\_T50 & 200 &  500 & 2 & 0.50 & gaussian & 0.75 & 0.8 \\
N200\_P1000 & 200 & 1000 & 1 & 0.50 & gaussian & 0.75 & 0.5 \\
N200\_P2000 & 200 & 2000 & 1 & 0.50 & gaussian & 0.75 & 0.5 \\
N400\_P1000 & 400 & 1000 & 1 & 0.50 & gaussian & 0.75 & 0.5 \\
N400\_P2000 & 400 & 2000 & 1 & 0.50 & gaussian & 0.75 & 0.5 \\
\bottomrule
\end{longtable}


\subsection{Monte Carlo uncertainty and selected targets}

For coordinate group $G$ in setting $a$, let
$C_{ar}=|G|^{-1}\sum_{k\in G}1\{\beta_k^\star\in I_{ark}\}$.
The setting's coverage is $B^{-1}\sum_r C_{ar}$ and its Monte Carlo
standard error is $\operatorname{sd}(C_{ar})/\sqrt B$.
A mean over $A$ settings has standard error
$A^{-1}\{\sum_a\widehat{\operatorname{se}}_a^2\}^{1/2}$.
The paired coverage difference applies this calculation to
$C_{ar}^{full}-C_{ar}^{ablated}$. It retains dependence between
coordinates and between procedures on the same dataset.

For an individual coordinate, exact binomial intervals use its
success count among 60 independent datasets.
Table~\ref{tab:joint-coordinate} also reports the mean estimated
standard error divided by the empirical standard deviation of that
coordinate's estimates. Empirical standard deviations are calculated
before combining different target coordinates.

Every selected Gaussian model has the population projection
\[
 \beta_{A,A}^P=\beta_A^0+
                    \Sigma_{AA}^{-1}\Sigma_{A,A^c}\beta_{A^c}^0.
\]
The saved selected sets and generating covariance determine its
coordinate displacement without additional Monte Carlo sampling.
The target comparison is performed separately in each estimation
fold and then combined with the estimator's fold weights. All
full-method fold targets agree with their full-design coordinate
at tolerance $10^{-8}$. The paired ablation has
\JointAblationShiftCount{} coordinate estimates with a larger
fold-averaged displacement, whose maximum absolute value is
\JointAblationMaxShift{}.

\begingroup\small
\begin{longtable}{lrrrrr}
\caption{Scenario results for the paired simulation comparison}\label{tab:joint-simulation-all}\\
\toprule
Setting & Coordinate & Full coverage & Ablation & Difference & Length \\
\midrule\endfirsthead
\toprule
Setting & Coordinate & Full coverage & Ablation & Difference & Length \\
\midrule\endhead
CORR\_Q2\_T50 & Active & 0.975 (0.014) & 0.983 (0.012) & -0.008 (0.015) & 0.489 / 0.472 \\
CORR\_Q2\_T50 & Null & 0.967 (0.016) & 0.975 (0.014) & -0.008 (0.015) & 0.568 / 0.545 \\
N200\_P1000 & Active & 0.942 (0.021) & 0.950 (0.020) & -0.008 (0.025) & 0.310 / 0.277 \\
N200\_P1000 & Null & 0.958 (0.022) & 0.958 (0.022) & 0.000 (0.012) & 0.343 / 0.311 \\
N200\_P2000 & Active & 0.975 (0.014) & 0.983 (0.012) & -0.008 (0.015) & 0.317 / 0.308 \\
N200\_P2000 & Null & 1.000 (0.000) & 0.958 (0.018) & 0.042 (0.018) & 0.371 / 0.338 \\
N200\_Q1\_T25 & Active & 0.967 (0.016) & 0.967 (0.016) & 0.000 (0.021) & 0.309 / 0.301 \\
N200\_Q1\_T25 & Null & 0.992 (0.008) & 0.975 (0.014) & 0.017 (0.017) & 0.395 / 0.327 \\
N200\_Q1\_T50 & Active & 0.958 (0.018) & 0.942 (0.024) & 0.017 (0.017) & 0.297 / 0.280 \\
N200\_Q1\_T50 & Null & 0.975 (0.014) & 0.967 (0.016) & 0.008 (0.015) & 0.330 / 0.317 \\
N200\_Q1\_T75 & Active & 0.950 (0.023) & 0.942 (0.021) & 0.008 (0.019) & 0.326 / 0.309 \\
N200\_Q1\_T75 & Null & 0.975 (0.014) & 0.958 (0.018) & 0.017 (0.020) & 0.386 / 0.342 \\
N200\_Q2\_T25 & Active & 0.975 (0.014) & 0.942 (0.024) & 0.033 (0.020) & 0.338 / 0.316 \\
N200\_Q2\_T25 & Null & 0.950 (0.020) & 0.967 (0.016) & -0.017 (0.017) & 0.362 / 0.333 \\
N200\_Q2\_T50 & Active & 0.983 (0.012) & 0.983 (0.012) & 0.000 (0.017) & 0.288 / 0.281 \\
N200\_Q2\_T50 & Null & 0.967 (0.016) & 0.958 (0.018) & 0.008 (0.015) & 0.347 / 0.309 \\
N200\_Q2\_T75 & Active & 0.975 (0.014) & 0.958 (0.018) & 0.017 (0.012) & 0.316 / 0.308 \\
N200\_Q2\_T75 & Null & 0.958 (0.022) & 0.958 (0.018) & 0.000 (0.012) & 0.385 / 0.335 \\
N400\_P1000 & Active & 0.958 (0.018) & 0.942 (0.027) & 0.017 (0.017) & 0.191 / 0.182 \\
N400\_P1000 & Null & 0.958 (0.022) & 0.950 (0.023) & 0.008 (0.015) & 0.219 / 0.199 \\
N400\_P2000 & Active & 0.958 (0.022) & 0.958 (0.018) & 0.000 (0.012) & 0.192 / 0.183 \\
N400\_P2000 & Null & 0.942 (0.021) & 0.933 (0.022) & 0.008 (0.019) & 0.214 / 0.197 \\
N400\_Q1\_T25 & Active & 1.000 (0.000) & 0.983 (0.012) & 0.017 (0.012) & 0.217 / 0.199 \\
N400\_Q1\_T25 & Null & 0.975 (0.014) & 0.950 (0.023) & 0.025 (0.019) & 0.235 / 0.223 \\
N400\_Q1\_T50 & Active & 0.967 (0.016) & 0.942 (0.024) & 0.025 (0.019) & 0.192 / 0.182 \\
N400\_Q1\_T50 & Null & 0.975 (0.014) & 0.967 (0.016) & 0.008 (0.015) & 0.213 / 0.197 \\
N400\_Q1\_T75 & Active & 0.942 (0.021) & 0.942 (0.021) & 0.000 (0.017) & 0.203 / 0.189 \\
N400\_Q1\_T75 & Null & 0.983 (0.012) & 0.975 (0.014) & 0.008 (0.015) & 0.221 / 0.207 \\
N400\_Q2\_T25 & Active & 0.975 (0.014) & 0.967 (0.016) & 0.008 (0.015) & 0.210 / 0.199 \\
N400\_Q2\_T25 & Null & 0.967 (0.016) & 0.958 (0.018) & 0.008 (0.008) & 0.230 / 0.218 \\
N400\_Q2\_T50 & Active & 0.958 (0.022) & 0.942 (0.027) & 0.017 (0.012) & 0.187 / 0.179 \\
N400\_Q2\_T50 & Null & 0.983 (0.012) & 0.950 (0.020) & 0.033 (0.016) & 0.206 / 0.198 \\
N400\_Q2\_T75 & Active & 0.942 (0.021) & 0.925 (0.023) & 0.017 (0.017) & 0.205 / 0.194 \\
N400\_Q2\_T75 & Null & 0.942 (0.021) & 0.967 (0.016) & -0.025 (0.022) & 0.226 / 0.213 \\
T3\_Q1\_T25 & Active & 0.967 (0.016) & 0.942 (0.021) & 0.025 (0.019) & 0.257 / 0.231 \\
T3\_Q1\_T25 & Null & 0.992 (0.008) & 0.967 (0.016) & 0.025 (0.019) & 0.301 / 0.261 \\
T3\_Q1\_T75 & Active & 0.967 (0.020) & 0.942 (0.024) & 0.025 (0.019) & 0.237 / 0.224 \\
T3\_Q1\_T75 & Null & 0.942 (0.021) & 0.958 (0.018) & -0.017 (0.020) & 0.299 / 0.264 \\
T3\_Q2\_T25 & Active & 0.967 (0.016) & 0.967 (0.016) & 0.000 (0.012) & 0.258 / 0.242 \\
T3\_Q2\_T25 & Null & 0.967 (0.016) & 0.967 (0.016) & 0.000 (0.012) & 0.281 / 0.247 \\
T3\_Q2\_T75 & Active & 0.950 (0.020) & 0.942 (0.021) & 0.008 (0.015) & 0.245 / 0.232 \\
T3\_Q2\_T75 & Null & 0.950 (0.020) & 0.950 (0.020) & 0.000 (0.021) & 0.281 / 0.258 \\
WEAK\_Q1\_T25 & Active & 0.942 (0.021) & 0.758 (0.044) & 0.183 (0.043) & 0.336 / 0.308 \\
WEAK\_Q1\_T25 & Null & 0.992 (0.008) & 0.942 (0.021) & 0.050 (0.023) & 0.398 / 0.366 \\
WEAK\_Q1\_T50 & Active & 0.983 (0.012) & 0.842 (0.037) & 0.142 (0.036) & 0.324 / 0.304 \\
WEAK\_Q1\_T50 & Null & 0.958 (0.018) & 0.917 (0.024) & 0.042 (0.027) & 0.361 / 0.326 \\
WEAK\_Q2\_T25 & Active & 0.950 (0.020) & 0.808 (0.036) & 0.142 (0.034) & 0.354 / 0.312 \\
WEAK\_Q2\_T25 & Null & 0.967 (0.016) & 0.908 (0.025) & 0.058 (0.021) & 0.400 / 0.333 \\
\bottomrule
\end{longtable}

\begin{longtable}{lrrrrr}
\caption{Coordinate coverage with exact binomial intervals and standard error ratios}\label{tab:joint-coordinate}\\
\toprule
Setting & Coordinate & Coverage & MCSE & Binomial interval & SE/SD \\
\midrule\endfirsthead
\toprule
Setting & Coordinate & Coverage & MCSE & Binomial interval & SE/SD \\
\midrule\endhead
CORR\_Q2\_T50 & x00001 & 0.967 & 0.023 & [0.885, 0.996] & 1.349 \\
CORR\_Q2\_T50 & x00002 & 0.983 & 0.017 & [0.911, 1.000] & 1.533 \\
CORR\_Q2\_T50 & x00006 & 0.933 & 0.032 & [0.838, 0.982] & 1.225 \\
CORR\_Q2\_T50 & x00007 & 1.000 & 0.000 & [0.940, 1.000] & 1.506 \\
N200\_P1000 & x00001 & 0.917 & 0.036 & [0.816, 0.972] & 1.112 \\
N200\_P1000 & x00002 & 0.967 & 0.023 & [0.885, 0.996] & 1.156 \\
N200\_P1000 & x00006 & 0.983 & 0.017 & [0.911, 1.000] & 1.327 \\
N200\_P1000 & x00007 & 0.933 & 0.032 & [0.838, 0.982] & 1.230 \\
N200\_P2000 & x00001 & 0.983 & 0.017 & [0.911, 1.000] & 1.399 \\
N200\_P2000 & x00002 & 0.967 & 0.023 & [0.885, 0.996] & 1.305 \\
N200\_P2000 & x00006 & 1.000 & 0.000 & [0.940, 1.000] & 1.418 \\
N200\_P2000 & x00007 & 1.000 & 0.000 & [0.940, 1.000] & 1.944 \\
N200\_Q1\_T25 & x00001 & 0.967 & 0.023 & [0.885, 0.996] & 1.284 \\
N200\_Q1\_T25 & x00002 & 0.967 & 0.023 & [0.885, 0.996] & 1.300 \\
N200\_Q1\_T25 & x00006 & 0.983 & 0.017 & [0.911, 1.000] & 1.467 \\
N200\_Q1\_T25 & x00007 & 1.000 & 0.000 & [0.940, 1.000] & 1.284 \\
N200\_Q1\_T50 & x00001 & 0.967 & 0.023 & [0.885, 0.996] & 1.255 \\
N200\_Q1\_T50 & x00002 & 0.950 & 0.028 & [0.861, 0.990] & 1.159 \\
N200\_Q1\_T50 & x00006 & 0.983 & 0.017 & [0.911, 1.000] & 1.254 \\
N200\_Q1\_T50 & x00007 & 0.967 & 0.023 & [0.885, 0.996] & 1.125 \\
N200\_Q1\_T75 & x00001 & 0.950 & 0.028 & [0.861, 0.990] & 1.269 \\
N200\_Q1\_T75 & x00002 & 0.950 & 0.028 & [0.861, 0.990] & 1.182 \\
N200\_Q1\_T75 & x00006 & 0.983 & 0.017 & [0.911, 1.000] & 1.259 \\
N200\_Q1\_T75 & x00007 & 0.967 & 0.023 & [0.885, 0.996] & 1.388 \\
N200\_Q2\_T25 & x00001 & 1.000 & 0.000 & [0.940, 1.000] & 1.250 \\
N200\_Q2\_T25 & x00002 & 0.950 & 0.028 & [0.861, 0.990] & 1.217 \\
N200\_Q2\_T25 & x00006 & 0.900 & 0.039 & [0.795, 0.962] & 1.118 \\
N200\_Q2\_T25 & x00007 & 1.000 & 0.000 & [0.940, 1.000] & 1.562 \\
N200\_Q2\_T50 & x00001 & 0.983 & 0.017 & [0.911, 1.000] & 1.249 \\
N200\_Q2\_T50 & x00002 & 0.983 & 0.017 & [0.911, 1.000] & 1.126 \\
N200\_Q2\_T50 & x00006 & 0.933 & 0.032 & [0.838, 0.982] & 1.129 \\
N200\_Q2\_T50 & x00007 & 1.000 & 0.000 & [0.940, 1.000] & 1.360 \\
N200\_Q2\_T75 & x00001 & 0.983 & 0.017 & [0.911, 1.000] & 1.197 \\
N200\_Q2\_T75 & x00002 & 0.967 & 0.023 & [0.885, 0.996] & 1.200 \\
N200\_Q2\_T75 & x00006 & 0.933 & 0.032 & [0.838, 0.982] & 1.275 \\
N200\_Q2\_T75 & x00007 & 0.983 & 0.017 & [0.911, 1.000] & 1.434 \\
N400\_P1000 & x00001 & 0.967 & 0.023 & [0.885, 0.996] & 1.129 \\
N400\_P1000 & x00002 & 0.950 & 0.028 & [0.861, 0.990] & 1.254 \\
N400\_P1000 & x00006 & 0.950 & 0.028 & [0.861, 0.990] & 1.227 \\
N400\_P1000 & x00007 & 0.967 & 0.023 & [0.885, 0.996] & 1.255 \\
N400\_P2000 & x00001 & 0.950 & 0.028 & [0.861, 0.990] & 1.079 \\
N400\_P2000 & x00002 & 0.967 & 0.023 & [0.885, 0.996] & 1.271 \\
N400\_P2000 & x00006 & 0.900 & 0.039 & [0.795, 0.962] & 0.977 \\
N400\_P2000 & x00007 & 0.983 & 0.017 & [0.911, 1.000] & 1.215 \\
N400\_Q1\_T25 & x00001 & 1.000 & 0.000 & [0.940, 1.000] & 1.429 \\
N400\_Q1\_T25 & x00002 & 1.000 & 0.000 & [0.940, 1.000] & 1.391 \\
N400\_Q1\_T25 & x00006 & 0.983 & 0.017 & [0.911, 1.000] & 1.242 \\
N400\_Q1\_T25 & x00007 & 0.967 & 0.023 & [0.885, 0.996] & 1.126 \\
N400\_Q1\_T50 & x00001 & 0.983 & 0.017 & [0.911, 1.000] & 1.196 \\
N400\_Q1\_T50 & x00002 & 0.950 & 0.028 & [0.861, 0.990] & 1.262 \\
N400\_Q1\_T50 & x00006 & 1.000 & 0.000 & [0.940, 1.000] & 1.191 \\
N400\_Q1\_T50 & x00007 & 0.950 & 0.028 & [0.861, 0.990] & 1.147 \\
N400\_Q1\_T75 & x00001 & 0.950 & 0.028 & [0.861, 0.990] & 1.155 \\
N400\_Q1\_T75 & x00002 & 0.933 & 0.032 & [0.838, 0.982] & 1.087 \\
N400\_Q1\_T75 & x00006 & 0.983 & 0.017 & [0.911, 1.000] & 1.131 \\
N400\_Q1\_T75 & x00007 & 0.983 & 0.017 & [0.911, 1.000] & 1.223 \\
N400\_Q2\_T25 & x00001 & 0.967 & 0.023 & [0.885, 0.996] & 1.063 \\
N400\_Q2\_T25 & x00002 & 0.983 & 0.017 & [0.911, 1.000] & 1.247 \\
N400\_Q2\_T25 & x00006 & 0.950 & 0.028 & [0.861, 0.990] & 1.098 \\
N400\_Q2\_T25 & x00007 & 0.983 & 0.017 & [0.911, 1.000] & 1.184 \\
N400\_Q2\_T50 & x00001 & 0.983 & 0.017 & [0.911, 1.000] & 1.113 \\
N400\_Q2\_T50 & x00002 & 0.933 & 0.032 & [0.838, 0.982] & 1.133 \\
N400\_Q2\_T50 & x00006 & 1.000 & 0.000 & [0.940, 1.000] & 1.302 \\
N400\_Q2\_T50 & x00007 & 0.967 & 0.023 & [0.885, 0.996] & 1.081 \\
N400\_Q2\_T75 & x00001 & 0.950 & 0.028 & [0.861, 0.990] & 1.072 \\
N400\_Q2\_T75 & x00002 & 0.933 & 0.032 & [0.838, 0.982] & 1.071 \\
N400\_Q2\_T75 & x00006 & 0.967 & 0.023 & [0.885, 0.996] & 1.097 \\
N400\_Q2\_T75 & x00007 & 0.917 & 0.036 & [0.816, 0.972] & 1.109 \\
T3\_Q1\_T25 & x00001 & 0.950 & 0.028 & [0.861, 0.990] & 1.121 \\
T3\_Q1\_T25 & x00002 & 0.983 & 0.017 & [0.911, 1.000] & 1.271 \\
T3\_Q1\_T25 & x00006 & 1.000 & 0.000 & [0.940, 1.000] & 1.581 \\
T3\_Q1\_T25 & x00007 & 0.983 & 0.017 & [0.911, 1.000] & 1.221 \\
T3\_Q1\_T75 & x00001 & 0.950 & 0.028 & [0.861, 0.990] & 1.013 \\
T3\_Q1\_T75 & x00002 & 0.983 & 0.017 & [0.911, 1.000] & 1.172 \\
T3\_Q1\_T75 & x00006 & 0.950 & 0.028 & [0.861, 0.990] & 1.170 \\
T3\_Q1\_T75 & x00007 & 0.933 & 0.032 & [0.838, 0.982] & 1.294 \\
T3\_Q2\_T25 & x00001 & 0.983 & 0.017 & [0.911, 1.000] & 1.332 \\
T3\_Q2\_T25 & x00002 & 0.950 & 0.028 & [0.861, 0.990] & 1.100 \\
T3\_Q2\_T25 & x00006 & 1.000 & 0.000 & [0.940, 1.000] & 1.460 \\
T3\_Q2\_T25 & x00007 & 0.933 & 0.032 & [0.838, 0.982] & 1.232 \\
T3\_Q2\_T75 & x00001 & 0.950 & 0.028 & [0.861, 0.990] & 1.222 \\
T3\_Q2\_T75 & x00002 & 0.950 & 0.028 & [0.861, 0.990] & 1.219 \\
T3\_Q2\_T75 & x00006 & 0.983 & 0.017 & [0.911, 1.000] & 1.157 \\
T3\_Q2\_T75 & x00007 & 0.917 & 0.036 & [0.816, 0.972] & 1.124 \\
WEAK\_Q1\_T25 & x00001 & 0.967 & 0.023 & [0.885, 0.996] & 1.120 \\
WEAK\_Q1\_T25 & x00002 & 0.917 & 0.036 & [0.816, 0.972] & 1.049 \\
WEAK\_Q1\_T25 & x00006 & 1.000 & 0.000 & [0.940, 1.000] & 1.388 \\
WEAK\_Q1\_T25 & x00007 & 0.983 & 0.017 & [0.911, 1.000] & 1.478 \\
WEAK\_Q1\_T50 & x00001 & 1.000 & 0.000 & [0.940, 1.000] & 1.253 \\
WEAK\_Q1\_T50 & x00002 & 0.967 & 0.023 & [0.885, 0.996] & 1.259 \\
WEAK\_Q1\_T50 & x00006 & 0.983 & 0.017 & [0.911, 1.000] & 1.302 \\
WEAK\_Q1\_T50 & x00007 & 0.933 & 0.032 & [0.838, 0.982] & 1.065 \\
WEAK\_Q2\_T25 & x00001 & 0.967 & 0.023 & [0.885, 0.996] & 1.392 \\
WEAK\_Q2\_T25 & x00002 & 0.933 & 0.032 & [0.838, 0.982] & 1.210 \\
WEAK\_Q2\_T25 & x00006 & 1.000 & 0.000 & [0.940, 1.000] & 1.293 \\
WEAK\_Q2\_T25 & x00007 & 0.933 & 0.032 & [0.838, 0.982] & 1.247 \\
\bottomrule
\end{longtable}

\endgroup

\subsection{Profile targets under predictor dependent scale}

The complete profile parameter can differ from both a generating
slope and a marginal quantile projection. A separate population
calculation uses innovation scale
$\exp\{0.35\operatorname{clip}(X_1,-3,3)\}$ in a Gaussian AR(1)
design with 1000 coordinates and eight nonzero coefficients.
Writing $X=aX_1+\eta$, where $a_k=\rho_X^{k-1}$ and $\eta$ is
independent of $X_1$, gives
\[
 E g_{i,h,k}(\beta^0+\delta e_1)
       =a_k E g_{i,h,1}(\beta^0+\delta e_1).
\]
The scalar population score therefore determines the full slope
adjustment. Four independent populations of $10^5$ clusters estimate
the complete profile target, with a small-bandwidth calculation and
a half-bandwidth refinement. The stored construction solves the
paired population score to approximate the common unsmoothed target.
Its role is target calculation; the sample estimator in this article
pairs fitted coefficients.

For the marginal projection
$(\alpha^{QR},\beta^{QR})\in\argmin_{\alpha,\beta}
E\rho_\tau(Y-\alpha-X^T\beta)$,
deterministic quadrature solves the intercept and first-slope
equations jointly. The other slopes equal their generating
coefficients by conditional Gaussian projection.
Table~\ref{tab:conditional-scale-targets} reports the two targets at
the Gaussian lower quartile and the $t_3$ upper quartile. Variation
between population repetitions gives the profile target's Monte Carlo
error. The target records retain the complete design, working
nuisance matrix and numerical residuals.

\begin{table}[t]
\caption{Profile and marginal quantile projections under predictor dependent error scales}\label{tab:conditional-scale-targets}
\begin{tabularx}{\textwidth}{@{}Xrrrrrr@{}}
\toprule
Error & Quantile & Generated & Profile & MCSE & QR slope & QR intercept \\
\midrule
Gaussian & 0.25 & 0.750 & 0.836 & 0.0009 & 0.817 & -0.282 \\
$t_3$ & 0.75 & 0.750 & 0.670 & 0.0013 & 0.693 & 0.406 \\
\bottomrule
\end{tabularx}
\par\smallskip\noindent Slope entries concern the first coordinate. Profile targets use four independent populations of 100{,}000 clusters and a half-bandwidth refinement; MCSE refers to that population calculation. Marginal QR includes an unpenalised intercept and uses deterministic quadrature, checked by four independent populations of 100{,}000 clusters. The fitted design has 1000 slope coordinates and eight nonzero generating coefficients.
\end{table}


\subsection{Gaussian influence diagnostics}
\newcommand{\GaussianDiagnosticDatasets}{10500}
\newcommand{\GaussianRemainderFiveSmall}{0.774}
\newcommand{\GaussianRemainderFiveLarge}{0.347}
\newcommand{\GaussianRemainderTenSmall}{1.159}
\newcommand{\GaussianRemainderTenLarge}{0.387}
\newcommand{\GaussianCurvatureIntegral}{0.3117484}
\newcommand{\GaussianCurvaturePopulation}{0.3119884}
\newcommand{\GaussianCurvatureMCSE}{0.0005868}
\newcommand{\GaussianCurvatureClusters}{100000}


The twelve Gaussian scaling settings retain all 10500 registered
datasets. The covariance is AR(1) with correlation $0.5$, cluster
sizes are uniform on three through eight, and the random-intercept
standard deviation is one. The median-quantile model has alternating
coefficients of magnitude $0.75$. Its population curvature is
given by \eqref{eq:supp-exact-gaussian-curvature}. Gaussian
quadrature gives $\kappa_0=\GaussianCurvatureIntegral$; an independent
population of \GaussianCurvatureClusters{} clusters gives
\GaussianCurvaturePopulation{} with Monte Carlo standard error
\GaussianCurvatureMCSE{}.

Each coordinate's remainder is
$\sqrt n\{\widehat\beta_k-\beta_k^0-n^{-1}\sum_i\xi_{ik,h}\}$.
Table~\ref{tab:gaussian-remainder} gives its root mean square and
Monte Carlo standard error. Within an active or null group, squared
remainders are averaged within each replication before aggregation.
The standard error uses the delta method at their square root.
The relative column replaces $\sqrt n$ by the inverse of that
replication's reported standard error before taking the root mean
square. The numerical records also retain fold-specific target
displacements and both bandwidth-specific Hessian errors.

\begingroup\small

\begin{longtable}[t]{lrrrlrrrr}
\caption{\label{tab:gaussian-remainder}Gaussian influence remainder by requested coordinate group}\\
\toprule
Setting & $n$ & $p$ & $s$ & Coordinates & Replications & RMS & MCSE & Relative\\
\midrule
A01 & 100 & 500 & 5 & Active & 1000 & 0.647 & 0.011 & 0.568\\
A01 & 100 & 500 & 5 & Null & 1000 & 0.758 & 0.013 & 0.569\\
A02 & 200 & 500 & 5 & Active & 1000 & 0.419 & 0.007 & 0.419\\
A02 & 200 & 500 & 5 & Null & 1000 & 0.490 & 0.008 & 0.432\\
A03 & 400 & 500 & 5 & Active & 1000 & 0.315 & 0.006 & 0.340\\
A03 & 400 & 500 & 5 & Null & 1000 & 0.382 & 0.007 & 0.364\\
A04 & 100 & 1000 & 5 & Active & 1000 & 0.703 & 0.012 & 0.617\\
A04 & 100 & 1000 & 5 & Null & 1000 & 0.839 & 0.015 & 0.616\\
A05 & 200 & 1000 & 5 & Active & 1000 & 0.418 & 0.007 & 0.414\\
A05 & 200 & 1000 & 5 & Null & 1000 & 0.493 & 0.008 & 0.423\\
A06 & 400 & 1000 & 5 & Active & 1000 & 0.312 & 0.006 & 0.333\\
A06 & 400 & 1000 & 5 & Null & 1000 & 0.380 & 0.006 & 0.364\\
A07 & 100 & 2000 & 5 & Active & 500 & 0.762 & 0.019 & 0.677\\
A07 & 100 & 2000 & 5 & Null & 500 & 0.908 & 0.022 & 0.653\\
A08 & 200 & 2000 & 5 & Active & 500 & 0.416 & 0.011 & 0.391\\
A08 & 200 & 2000 & 5 & Null & 500 & 0.510 & 0.013 & 0.425\\
A09 & 400 & 2000 & 5 & Active & 500 & 0.311 & 0.007 & 0.330\\
A09 & 400 & 2000 & 5 & Null & 500 & 0.384 & 0.010 & 0.362\\
A10 & 100 & 1000 & 10 & Active & 1000 & 1.087 & 0.018 & 0.791\\
A10 & 100 & 1000 & 10 & Null & 1000 & 1.226 & 0.021 & 0.762\\
A11 & 200 & 1000 & 10 & Active & 1000 & 0.532 & 0.010 & 0.451\\
A11 & 200 & 1000 & 10 & Null & 1000 & 0.585 & 0.010 & 0.439\\
A12 & 400 & 1000 & 10 & Active & 1000 & 0.355 & 0.006 & 0.358\\
A12 & 400 & 1000 & 10 & Null & 1000 & 0.416 & 0.007 & 0.379\\
\bottomrule
\end{longtable}

\endgroup

\subsection{Runtime and numerical records}

The scaling experiment uses 20 timed fits at each of
$p=500,1000,2000,10000$, with $n=200$, $q=1$, $\tau=0.50$ and
four requested coordinates. Data generation is excluded from the
fitting clock. Selection and inference times are retained separately,
together with the largest selected model and final score residual.
The executable source and saved session information identify the
machine and numerical environment for these measurements.

\newcommand{\SubmissionTimingReplications}{350}
\newcommand{\SubmissionTimingJointHigh}{3.031}
\newcommand{\SubmissionTimingPooledHigh}{80.340}
\newcommand{\SubmissionTimingTwelve}{4.471}
\newcommand{\SubmissionWallFour}{331.943}
\newcommand{\SubmissionWallEight}{197.697}

The registered computation study contains
\SubmissionTimingReplications{} timed replications in the seven
settings of Table~\ref{tab:registered-computation-design}. The
available-predictor column gives the pool passed to each method
after any prescreening. The proposed procedure subsequently selects
its coordinate-specific designs. In G04, an outcome-free variance
filter retains 500 of the original 5000 predictors, together with
any requested coordinates outside that set. The known-support oracle
uses the generated active coordinates and the requested targets.

\begin{table}

\caption{\label{tab:registered-computation-design}Registered computation settings with fifty timed repetitions}
\centering
\begin{tabular}[t]{lrrrrrr}
\toprule
Setting & $n$ & $p$ & Fitted & $q$ & Targets & Workers\\
\midrule
G01 & 100 & 500 & 500 & 1 & 4 & 1\\
G02 & 200 & 2000 & 2000 & 1 & 4 & 1\\
G03 & 200 & 1000 & 1000 & 2 & 4 & 1\\
G04 & 400 & 5000 & 500 & 1 & 4 & 1\\
G05 & 200 & 1000 & 1000 & 1 & 12 & 1\\
G06 & 200 & 1000 & 1000 & 1 & 4 & 4\\
G07 & 200 & 1000 & 1000 & 1 & 4 & 8\\
\bottomrule
\end{tabular}
\end{table}


Table~\ref{tab:registered-computation-time} reports median fitting
times and empirical interquartile ranges over fifty runs. These
clocks include tuning and estimation after prescreening. The complete
batch clocks also include data generation, prescreening, process
startup and result storage. During timing, the other simulation and
application fitting processes in this project were suspended.
All registered method fits converged.

G06 and G07 use the same fifty generated datasets and the same
method-specific tuning seeds. Four workers complete their batches
in \SubmissionWallFour{} seconds and eight workers in
\SubmissionWallEight{} seconds. The paired method comparisons return
identical coefficient vectors, interval endpoints and standard
errors. Each batch includes the proposed method, pooled quantile
Lasso and the known-support oracle.

\begingroup\small

\begin{longtable}[t]{lllll}
\caption{\label{tab:registered-computation-time}Elapsed fitting time in seconds including each method's tuning}\\
\toprule
Setting & Method & Fits & Median & Quartiles\\
\midrule
G01 & Joint profile & 50/50 & 0.663 & {}[0.620, 0.695]\\
G01 & Known support & 50/50 & 0.004 & {}[0.004, 0.004]\\
G01 & Pooled Lasso & 50/50 & 2.572 & {}[2.419, 2.693]\\
G02 & Joint profile & 50/50 & 3.031 & {}[2.944, 3.133]\\
G02 & Known support & 50/50 & 0.010 & {}[0.010, 0.012]\\
G02 & Pooled Lasso & 50/50 & 80.340 & {}[78.450, 84.394]\\
G03 & Joint profile & 50/50 & 1.781 & {}[1.702, 1.849]\\
G03 & Known support & 50/50 & 0.009 & {}[0.008, 0.010]\\
G03 & Pooled Lasso & 50/50 & 20.113 & {}[19.157, 20.892]\\
G04 & Joint profile & 50/50 & 2.709 & {}[2.600, 2.866]\\
G04 & Known support & 50/50 & 0.033 & {}[0.032, 0.037]\\
G04 & Pooled Lasso & 50/50 & 10.977 & {}[10.644, 11.469]\\
G05 & Joint profile & 50/50 & 4.471 & {}[4.405, 4.583]\\
G05 & Known support & 50/50 & 0.009 & {}[0.008, 0.009]\\
G06 & Joint profile & 50/50 & 2.046 & {}[1.983, 2.111]\\
G06 & Known support & 50/50 & 0.009 & {}[0.008, 0.014]\\
G06 & Pooled Lasso & 50/50 & 22.698 & {}[21.894, 23.623]\\
G07 & Joint profile & 50/50 & 2.310 & {}[2.215, 2.400]\\
G07 & Known support & 50/50 & 0.013 & {}[0.011, 0.028]\\
G07 & Pooled Lasso & 50/50 & 26.155 & {}[25.188, 27.128]\\
\bottomrule
\end{longtable}

\endgroup

The independent evaluation and its ablation each retain a result
for every registered dataset and coordinate. The final result
builder requires complete paired keys before calculating summaries.
It generates the tables from saved numerical records, including
interval lengths and Monte Carlo errors, without manual insertion
of result values.
